\section{Introduction}
\label{sec:intro}

Look at the worker in \cref{fig:teaser}. He wears a hard hat, holds a cigarette between his lips, and has neither gloves nor a mask. Under the inspection rule of a coal-conveyor gallery, smoking is a level-one hazard, so the scene is level one. A vision-language model (VLM) fine-tuned on such verdicts gets it right. Now the regulation is revised and smoking becomes a level-two hazard. The scene should drop to level two---but the fine-tuned model insists on level one, because it ``sees'' an open flame with probability 0.95. There is no flame in the image. The untrained model, asked the same question, does not see one either.

This is not an isolated glitch, and it is not a lack of capacity. It is a consequence of how verdicts supervise facts. Under the original rule, smoking alone decides the verdict, so the open-flame fact can take any value without changing the label: the verdict carries \emph{no information} about it. In the training set, every one of the 71 images with smoking comes from an industry where smoking is level one, so not a single training label ever said ``this is a cigarette, not a flame''. We call such facts \emph{overexposed}. As with the highlights of an overexposed photograph, the detail is gone while the picture still looks right, and the loss becomes visible only when the exposure changes---here, when the rule is edited.

Rule-governed visual decisions are everywhere: hazard levels in industrial inspection, accept/reject in quality control, stop/go in driving. Their rules are written down, and they change: regulations are revised, sites have their own severity tables, an emergency suspends a traffic rule. A model that has to be retrained for every revision is impractical, and a model that silently breaks after one is dangerous.

\paragraph{Overexposure, quantified.} We make the observation precise (\cref{sec:theory}). The \emph{exposure} of a fact is the probability that flipping it changes the verdict. We show that the information a verdict carries about a fact is exactly its exposure times its entropy, and that the error after a rule edit is bounded by each fact's exposure under the \emph{new} rule. Facts with zero exposure under the training rule but positive exposure under a future rule are unidentifiable from verdicts---for any loss and any optimiser---yet they decide the edited outcome. The theory also predicts \emph{when} overexposure bites: when a rule is relaxed, because relaxing a fact exposes the facts it used to clip.

\paragraph{\lux.} The analysis suggests learning the state of the scene rather than the verdict of one rule (\cref{sec:method}). \lux asks a fixed set of atomic yes/no questions in parallel over one shared image prefill---at the latency of asking for the verdict directly---and composes the verdict with an editable severity table. Its training objective follows from three results: the posterior over facts is the only object that is optimal for every rule at once; observed facts are fitted through an explicit model of how annotations miss them; and facts the data cannot identify are kept at the base model (minimum-information completion).

\paragraph{Evidence.} On three datasets---industrial safety inspection (InspecSafe), industrial defects (MMAD) and driving (BDD-OIA)---pre-registered interventions show the predicted signature: verdict training raises false positives on clipped facts by 1.7--4.6$\times$ while visible facts are unchanged (\cref{sec:exp-over}). On real regulation edits scored against blind labels from three annotators, \lux is correct on 89\% of the affected scenes; a fine-tuned model that is told the new rule is correct on 26\%, and even one retrained on relabelled data on 57\% (\cref{sec:exp-edit}). With complete fact annotations the repair costs nothing in-distribution; with incomplete annotations the cost grows predictably with the missing rate, and modelling the missing annotations is several times cheaper than treating them as negatives (\cref{sec:exp-cost}).

\paragraph{Contributions.}
\begin{itemize}
\item We identify \emph{overexposure}, the loss of facts a training rule does not need, and characterise it with an information-theoretic measure (exposure) and a regret bound under rule edits.
\item We derive \lux, a rule-editable VLM reader whose training objective follows from the analysis rather than from heuristics.
\item We provide pre-registered evidence on three datasets, including blind human relabelling of real regulation edits, and report every hypothesis we tested, including the failed ones.
\end{itemize}

\paragraph{What is not new.} The phenomenon has relatives in many fields: explaining away~\cite{pearl1988}, the discriminative-part bias of weakly supervised localisation~\cite{zhou2016cam,singh2017hide}, gradient starvation~\cite{pezeshki2021gradient}, concept leakage~\cite{mahinpei2021promises,havasi2022leakage} and reasoning shortcuts~\cite{marconato2023not,bortolotti2025shortcuts}, whose benchmark already uses an edited traffic rule as a probe on synthetic scenes~\cite{bortolotti2024rsbench}. What we add is a graded measure with a deployment bound, a repair that needs no side channel, and evidence on real images and real rules.
